\subsection{Logical foundations}
\label{sec:logical-foundations}
\slidesrc{Prolog III: 1--9}

Why does Prolog compute correct answers? Because running a Prolog program is building a
\emph{proof} in first-order logic, with a single, very efficient inference rule:
\textbf{SLD resolution}. This deck was shown last; it ties Part~IV back to Part~II.

\subsubsection{Clauses and Horn clauses}
\slidesrc{Prolog III: 4--6}

\begin{nbdef}{Clause (slide 4)}
A \textbf{clause} is a disjunction of atoms or their negations (\textbf{literals}), where the
variables are implicitly \emph{universally quantified}:
\[
A_1 \lor A_2 \lor \cdots \lor A_n \lor \neg B_1 \lor \neg B_2 \lor \cdots \lor \neg B_m
\]
It is equivalent to
\[
(A_1 \lor \cdots \lor A_n) \leftarrow (B_1 \land \cdots \land B_m)
\]
usually written (conclusion $\leftarrow$ precondition) as
$A_1, \ldots, A_n \leftarrow B_1, \ldots, B_m$.
\end{nbdef}

Why are the two forms equivalent? With the rules of Section~\ref{sec:tautologies}:
\begin{align*}
&(A_1 \lor \cdots \lor A_n) \leftarrow (B_1 \land \cdots \land B_m) \\
&\quad\equiv \neg(B_1 \land \cdots \land B_m) \lor (A_1 \lor \cdots \lor A_n)
     && \reason{$q \leftarrow p \;\equiv\; \neg p \lor q$} \\
&\quad\equiv (\neg B_1 \lor \cdots \lor \neg B_m) \lor A_1 \lor \cdots \lor A_n
     && \reason{De Morgan} \\
&\quad\equiv A_1 \lor \cdots \lor A_n \lor \neg B_1 \lor \cdots \lor \neg B_m
     && \reason{commutativity}
\end{align*}
The positive literals end up on the left (conclusions), the negative ones on the right
(conditions).

\begin{nbdef}{Horn clause (slide 5)}
A \textbf{Horn clause} is a clause with \textbf{at most one positive literal} (a literal is a
predicate symbol applied to its arguments, or its negation).
\begin{center}
\small
\begin{tabular}{@{}l l l l@{}}
\toprule
\textbf{Kind} & \textbf{Logic} & \textbf{Prolog} & \textbf{Positive literals} \\
\midrule
definite clause: \textbf{rule} & $A \leftarrow B_1, \ldots, B_m$ & \code{a :- b1, ..., bm.} & 1 \\
definite clause: \textbf{fact} & $A \leftarrow$ & \code{a.} & 1 \\
negative clause: \textbf{goal} & $\leftarrow B_1, \ldots, B_m$ & \code{?- b1, ..., bm.} & 0 \\
\bottomrule
\end{tabular}
\end{center}
\end{nbdef}

\begin{nbfact}{Programs and questions (slide 6)}
\begin{itemize}[nosep]
  \item \textbf{Definite clauses represent knowledge}: a \textbf{logic program} is a set of
        definite clauses.
  \item \textbf{Negative clauses are questions}: the query \code{?- b1, ..., bm.} asks
        whether $\exists X_1 \ldots X_n.\,(B_1 \land \cdots \land B_m)$ holds (the
        $X_i$ are the variables of the query).
  \item To answer, Prolog uses \textbf{refutation}: it adds the \emph{negation} of the
        question and tries to derive a contradiction:
        \begin{align*}
        &\neg\,\exists X_1 \ldots X_n.\,(B_1 \land \cdots \land B_m) \\
        &\quad\equiv \forall X_1 \ldots X_n.\,(\neg B_1 \lor \cdots \lor \neg B_m)
          && \reason{duality, De Morgan} \\
        &\quad\equiv\; \leftarrow B_1, \ldots, B_m
          && \reason{a negative clause}
        \end{align*}
\end{itemize}
\end{nbfact}

\begin{nbintuition}[Refutation = proof by contradiction]
It is the $\neg$\,Intro idea of natural deduction (Section~\ref{sec:natural-deduction-prop}):
to show that ``someone is mortal'', assume ``nobody is mortal'' and derive a contradiction
from it and the program. The contradiction is the \textbf{empty clause} $\Box$ (a disjunction
with no literals, which can never be true). The bindings used along the way tell \emph{who}
is mortal: they are the answer.
\end{nbintuition}

\subsubsection{SLD resolution}
\slidesrc{Prolog III: 7}

\begin{nbdef}{SLD resolution (slide 7)}
\textbf{S}election-function, \textbf{L}inear resolution for \textbf{D}efinite clauses.
Start from a definite program $P$ and a goal
\[ G_0 = \;\leftarrow A_1, \ldots, A_m \qquad (m \geq 0). \]
\begin{enumerate}[nosep]
  \item The \textbf{selection rule} selects a sub-goal $A_i$.
  \item Choose a program clause, \textbf{renamed} with fresh variables,
        $B_0 \leftarrow B_1, \ldots, B_n$ ($n \geq 0$), whose head $B_0$ unifies with $A_i$,
        with most general unifier $\theta_1$.
  \item The new goal is
        \[ G_1 = \;\leftarrow (A_1, \ldots, A_{i-1}, B_1, \ldots, B_n, A_{i+1}, \ldots, A_m)\,\theta_1 \]
        i.e.\ $A_i$ is \emph{replaced by the body} of the clause, and the unifier is applied
        to the whole goal.
\end{enumerate}
Repeat until the goal is empty ($\Box$: success) or no clause unifies with the selected atom
(failure of that branch). The answer is the composition of all the unifiers, restricted to
the variables of $G_0$.
\end{nbdef}

\begin{nbexample}{SLD derivation of \texttt{?- mortal(X).}}
Program: \code{mortal(X) :- human(X).} \ \code{human(socrates).}
\begin{align*}
G_0 &= \;\leftarrow \mathit{mortal}(X) \\
  &\quad \text{clause renamed } \mathit{mortal}(X_1) \leftarrow \mathit{human}(X_1),\quad
     \theta_1 = \{X_1 / X\} \\
G_1 &= \;\leftarrow \mathit{human}(X) \\
  &\quad \text{clause } \mathit{human}(\mathit{socrates}) \leftarrow,\quad
     \theta_2 = \{X / \mathit{socrates}\} \\
G_2 &= \;\Box \qquad \text{(empty goal: success)}
\end{align*}
Answer: $X = \mathit{socrates}$. It is Whiteboard~\ref{wb:modus-ponens} again: unification
does the instantiation, resolution does the Modus Ponens.
\end{nbexample}

\subsubsection{SLD trees}
\slidesrc{Prolog III: 8--9}

\begin{nbdef}{SLD tree (slide 8)}
Let $P$ be a definite program, $G_0$ a goal and $\mathcal{R}$ a selection rule. The SLD tree
of $G_0$ (using $P$ and $\mathcal{R}$) is a (possibly infinite) labelled tree such that:
\begin{itemize}[nosep]
  \item the root is labelled by $G_0$;
  \item if $\leftarrow A_1, \ldots, A_n$ is a node and $A_m$ is the atom selected by
        $\mathcal{R}$, the node has one child for each clause $A \leftarrow B_1, \ldots, B_q$
        of $P$ such that $A$ and $A_m$ are unifiable; the child is
        \[ \leftarrow (A_1, \ldots, A_{m-1}, B_1, \ldots, B_q, A_{m+1}, \ldots, A_n)\,\theta \]
        where $\theta$ is the mgu of $A_m$ and $A$.
\end{itemize}
Leaves labelled $\Box$ are \textbf{success} leaves (each gives an answer); leaves where no
clause applies are \textbf{failure} leaves.
\end{nbdef}

(The slide writes the child as $(A_1, \ldots, A_{m-1}, B_1, \ldots, A_{m+1}, \ldots, A_n)\theta$,
dropping $B_q$: the whole body $B_1, \ldots, B_q$ replaces $A_m$.)

\begin{nbexample}{SLD tree with the \emph{rightmost} selection rule (slide 9)}
Program (colours as on the slide):
\begin{center}
\textcolor{nbred}{(1)\ $p(X, Z) \leftarrow p(Y, Z), q(X, Y)$} \qquad
\textcolor{nbblue}{(2)\ $p(X, X)$} \qquad
\textcolor{nbgreen!60!black}{(3)\ $q(a, b)$}
\end{center}
Goal $\leftarrow p(X, b)$; the selection function picks the \textbf{rightmost} atom.
\begin{center}
\begin{forest}
for tree={draw, rounded corners, font=\small, l sep=8mm, s sep=10mm, edge={thick}, align=center}
[{$\leftarrow p(X, b)$}, fill=nbblue!8
  [{$\leftarrow p(Y, b),\ q(X, Y)$}, fill=nbblue!8, edge={nbred}, edge label={node[midway, left, font=\scriptsize]{(1)}}
    [{$\leftarrow p(b, b)$}, fill=nbblue!8, edge={nbgreen!60!black}, edge label={node[midway, right, font=\scriptsize]{(3): $X = a,\ Y = b$}}
      [{$\leftarrow p(Y', b),\ q(b, Y')$\\failure: no $q(b, \_)$}, fill=nbred!10, edge={nbred}, edge label={node[midway, left, font=\scriptsize]{(1)}}]
      [{$\Box$\\answer $X = a$}, fill=nbgreen!15, edge={nbblue}, edge label={node[midway, right, font=\scriptsize]{(2)}}]]]
  [{$\Box$\\answer $X = b$}, fill=nbgreen!15, edge={nbblue}, edge label={node[midway, right, font=\scriptsize]{(2): $X = b$}}]]
\end{forest}
\end{center}
Step by step:
\begin{enumerate}[nosep]
  \item Root $\leftarrow p(X, b)$: only one atom. Clause (1), renamed
        $p(X_1, Z_1) \leftarrow p(Y_1, Z_1), q(X_1, Y_1)$, unifies with $X_1 = X$, $Z_1 = b$:
        child $\leftarrow p(Y, b), q(X, Y)$ (writing $Y$ for $Y_1$). Clause (2), $p(X_2,
        X_2)$, unifies with $X_2 = X = b$: child $\Box$, answer $X = b$.
  \item Node $\leftarrow p(Y, b), q(X, Y)$: rightmost atom $q(X, Y)$; only clause (3)
        unifies: $X = a$, $Y = b$; child $\leftarrow p(b, b)$.
  \item Node $\leftarrow p(b, b)$: clause (1) gives $\leftarrow p(Y', b), q(b, Y')$, whose
        rightmost atom $q(b, Y')$ unifies with no clause: failure. Clause (2) unifies
        $p(b, b)$ with $p(X_2, X_2)$: $\Box$, answer $X = a$.
\end{enumerate}
The tree is \emph{finite}: two answers, $X = b$ and $X = a$.
\end{nbexample}

\begin{nbwarning}[What real Prolog does with the same program]
\notinslides
Prolog's selection rule is \textbf{leftmost}, its search is \textbf{depth-first} with clauses
tried \textbf{top to bottom}. With the leftmost rule the first child of
$\leftarrow p(X, b)$ is $\leftarrow p(Y, b), q(X, Y)$, whose leftmost atom $p(Y, b)$ is again
resolved first with clause (1): $\leftarrow p(Y', b), q(Y, Y'), q(X, Y)$, and so on. The
leftmost branch of the SLD tree is \emph{infinite}, and depth-first search explores it first:
\code{?- p(X, b).} never returns an answer (stack overflow), although $X = b$ and $X = a$ are
logical consequences of the program.
\begin{itemize}[nosep]
  \item The \emph{set of answers} does not depend on the selection rule (every SLD tree
        contains the same success leaves), but the \emph{shape} of the tree, and hence
        termination, does.
  \item Depth-first search is fast and memory-cheap, but \emph{incomplete}: it can get lost in
        an infinite branch before reaching a success leaf. This is the ``left recursion''
        problem of Section~\ref{sec:prolog-basics}. Writing clause (2) first, or the body as
        \code{q(X, Y), p(Y, Z)}, would make Prolog find both answers.
\end{itemize}
\end{nbwarning}

\begin{center}
\begin{tikzpicture}[font=\small, >=Stealth,
  b/.style={draw, thick, rounded corners, text width=3.6cm, align=center, minimum height=1.3cm, fill=#1}]
  \node[b=nbblue!8] (fol) at (0,0) {\textbf{FOL}\\full first-order logic, undecidable};
  \node[b=nbpurple!8] (horn) at (4.5,0) {\textbf{Horn clauses}\\facts, rules, goals};
  \node[b=nbgreen!10] (sld) at (9,0) {\textbf{SLD resolution}\\+ leftmost, depth-first\\= \textbf{Prolog}};
  \draw[->, thick] (fol) -- node[above, font=\scriptsize]{restrict} (horn);
  \draw[->, thick] (horn) -- node[above, font=\scriptsize]{one rule} (sld);
\end{tikzpicture}
\captionof{figure}{From logic to Prolog: restrict FOL to Horn clauses, use one inference rule
(SLD resolution), and fix a search strategy. The price of efficiency is that some logically
valid answers may not be found (non-termination).}
\end{center}
